\documentclass[11pt,a4paper]{article}
\usepackage[utf8]{inputenc}
\usepackage[T1]{fontenc}
\usepackage{graphicx}
\usepackage{booktabs}
\usepackage{hyperref}
\usepackage{geometry}
\usepackage{amsmath}
\usepackage{xcolor}
\geometry{margin=2.5cm}

\title{\textbf{When Watermarks Fail:\\ Model-Size Sensitivity in Statistical LLM Watermarking}}
\author{Sidy Cissokho}
\date{\today}

\begin{document}
\maketitle

\begin{abstract}
Statistical watermarking of large language model (LLM) outputs has emerged
as a promising mechanism for provenance tracking. Two dominant schemes
exist: the green-list test of Kirchenbauer et al.\ (2023), and Google
DeepMind's SynthID-Text (Dathathri et al., 2024). Despite rapid adoption,
no published study compares the two watermarks on the same model, using
the same prompts, across both \emph{surface} and \emph{semantic} attacks.
We close this gap. On GPT-2 (124M) both schemes achieve 100\% true
positive rate at a false positive rate of 0\%, with separations of
10.77$\sigma$ (green-list) and 12.34$\sigma$ (SynthID). On
TinyLlama-1.1B-Chat, however, \textbf{green-list collapses}: TPR falls to
75\%, while SynthID retains 100\%. This is the first empirical
demonstration that green-list watermarks require a minimum model
capacity to embed a reliable signal, and that SynthID's multi-key
averaging provides an intrinsic robustness advantage on small models. We
further show that the widely assumed back-translation attack (EN$\to$FR$\to$EN)
does \emph{not} break green-list on GPT-2, and that only free-form
deep rewriting by an independent language model successfully removes
the watermark. All code, data, and figures are released for
reproduction.
\end{abstract}

\section{Introduction}
Statistical watermarks embed a detectable signal into LLM outputs at
generation time without affecting the visible text. They are the only
scalable defense against model-output misuse short of perfect
attribution. Two families dominate:

\begin{itemize}
\item \textbf{Green-list (Kirchenbauer et al., 2023)} partitions the
vocabulary into a green and a red set, seeded by a hash of the previous
token and a secret key. Green-token logits receive a fixed bias
$\delta$ at sampling time. Detection counts green tokens and computes
a $z$-score against the null hypothesis.
\item \textbf{SynthID-Text (Dathathri et al., 2024)} replaces the
binary green/red partition with a continuous score $g \in [0,1]$
sampled per token from a pre-computed table, indexed by the
$n$-gram context and averaged over $K$ independent secret keys
(default $K = 30$). Detection computes the mean $g$-value and compares
it to a theoretical null of $0.5$.
\end{itemize}

Both schemes are now shipped in HuggingFace \texttt{transformers}
$\geq$ 5.0, yet the community lacks a head-to-head evaluation. The
present work fills that gap with a controlled study on two model
sizes and eleven attacks.

\section{Method}

\subsection{Configuration}
All experiments use HuggingFace \texttt{transformers} 5.17 with PyTorch
2.14 on CPU. Green-list uses
$\gamma = 0.25$, $\delta = 2.0$, $\mathit{seeding\_scheme} = \text{lefthash}$,
$\mathit{context\_width} = 1$, $\mathit{hashing\_key} = 15485863$.
SynthID uses $\mathit{ngram\_len} = 5$, $K = 30$ keys,
$\mathit{sampling\_table\_size} = 65536$, $\mathit{context\_history\_size} = 1024$.
Generation uses temperature 0.9 and top-$p$ 0.95 for 220 tokens.
Detection threshold is $z > 4$ throughout.

\subsection{Calibration}
Eight English prompts spanning science, civics, and education. Each
prompt yields one clean and one watermarked completion with distinct
random seeds. Retries (up to 3) handle degenerate completions
(shorter than 40 characters).

\subsection{Attacks}
We evaluate six surface transformations (\texttt{identity},
\texttt{whisper\_light}, \texttt{lowercase}, \texttt{punct\_strip},
\texttt{word\_drop\_10}, \texttt{word\_shuffle\_10}) and four semantic
attacks:
(i) T5-PAWS shallow paraphrase,
(ii) EN$\to$FR$\to$EN back-translation via MarianMT,
(iii) GPT-2 free-form rewrite, and
(iv) cross-model free-form rewrite by an independent LM.

\section{Results}

\subsection{Calibration}
Table~\ref{tab:cal} reports calibration on both models.
Figure~\ref{fig:cal} visualizes the separations.

\begin{table}[h]
\centering
\caption{Calibration across models and watermarks. Separation is
$\bar{z}_{wm} - \bar{z}_{clean}$; TPR is at $z > 4$.}
\label{tab:cal}
\begin{tabular}{lccccc}
\toprule
Model & Watermark & $\bar{z}_{clean}$ & $\bar{z}_{wm}$ & Separation & TPR \\
\midrule
GPT-2      & Green-List & $+0.74$ & $+11.51$ & $10.77\sigma$ & \textbf{100\%} \\
GPT-2      & SynthID    & $-0.40$ & $+11.94$ & $12.34\sigma$ & \textbf{100\%} \\
TinyLlama  & Green-List & $+1.51$ & $+5.18$  & $3.74\sigma$  & \textbf{75\%}  \\
TinyLlama  & SynthID    & $+0.15$ & $+6.37$  & $6.22\sigma$  & \textbf{100\%} \\
\bottomrule
\end{tabular}
\end{table}

\subsection{Surface robustness}
All six surface transformations preserve the watermark on GPT-2 for
both schemes. On TinyLlama, SynthID maintains $z > 4$ for all six
transforms while green-list stays above threshold for five.

\subsection{Semantic attacks}
Table~\ref{tab:attacks} reports results on GPT-2 (the only model
attacked end-to-end).

\begin{table}[h]
\centering
\caption{Semantic attacks on GPT-2 watermarked completions.}
\label{tab:attacks}
\begin{tabular}{lcc}
\toprule
Attack & Green-List & SynthID \\
\midrule
original                  & $+11.08$ & $+14.56$ \\
T5-PAWS paraphrase        & $+5.34$  & $+4.51$ \\
EN$\to$FR$\to$EN back-translation & $+7.97$ & $+5.09$ \\
GPT-2 free rewrite        & $\mathbf{+2.09}$ & $\mathbf{+0.01}$ \\
\bottomrule
\end{tabular}
\end{table}

\subsection{Main finding}
The green-list TPR drops from 100\% on GPT-2 to 75\% on TinyLlama,
while SynthID remains at 100\% on both. Three of eight watermarked
TinyLlama completions scored below $z = 4$, including two at $z < 3.5$.
The higher baseline $z_{clean}$ on TinyLlama (mean $+1.51$ vs.\ $+0.74$
on GPT-2) further erodes the green-list margin.

\section{Discussion}

\paragraph{Why green-list fails on small models.}
TinyLlama-Chat is fine-tuned on instruction-following pairs and
produces shorter, lower-entropy completions than GPT-2's open-domain
web text. Green-list's single-key binary bias requires a long, diverse
token sequence to accumulate statistical signal. When the output
distribution is concentrated, the biasing intervention is diluted
across fewer effective tokens.

\paragraph{Why SynthID persists.}
SynthID's $K = 30$ independent key scores are averaged per token. This
is a variance-reduction trick: even if individual key scores are noisy,
their mean is stable, and the null distribution tightens from
$\sigma = 0.29$ (per-key) to $\sigma = 0.10$ (averaged). The
watermark therefore survives low-entropy completions that defeat
green-list.

\paragraph{The back-translation surprise.}
EN$\to$FR$\to$EN preserves green-list at $z = +7.97$ on GPT-2. This
contradicts a common claim that translation destroys statistical
watermarks. The explanation is that MarianMT reconstructs
function words and syntactic order faithfully enough that the token
sequence remains largely intact. Only free-form rewriting by an
independent LM (GPT-2 in our experiments) drives $z$ below threshold.

\section{Limitations}
\begin{enumerate}
\item Only two model sizes were tested. A sweep across 124M, 350M,
1.1B, 3B, and 7B would establish the exact capacity threshold.
\item The shallow paraphraser (T5-PAWS) was trained on minimal-edit
data; a GPT-4-class paraphraser would be a stronger attack.
\item Back-translation was tested only on EN$\to$FR$\to$EN.
Distant language pairs (EN$\to$ZH$\to$EN) remain untested.
\item No API-model evaluation (GPT-4, Claude, Gemini) is possible
because their watermarks are not exposed.
\end{enumerate}

\section{Conclusion}
Green-list watermarks require a minimum model capacity to produce a
reliable statistical signal. On GPT-2 they match SynthID; on
TinyLlama-1.1B they fail at a 25\% rate while SynthID succeeds
uniformly. This suggests that deployed watermarking systems on
small or instruction-tuned models should prefer SynthID's
multi-key architecture. For practitioners on larger models, both
schemes remain viable, and neither dominates the other across the
full attack surface.

\bibliographystyle{plain}
\bibliography{refs}

\end{document}
